% -*-coding: utf-8 -*-
% File: preface/cover.tex
%
% 修改自: PlutoThesis_UTF8_1.9.4.20100419.zip
%         http://code.google.com/p/plutothesis/
%
%
%% 封面和摘要


%% 封面内容

% 授权书中用到的论文题目, 无需手动断行
% 使用该命令的文件:
%    appendix/Authorization.tex

\newcommand{\chinesethesistitle}{水下绳驱机械臂结构设计及研制}



% 英文论文题目
% Tip: \uppercase 作用是将英文标题字母全部大写

\newcommand{\englishthesistitle}{\uppercase{\LaTeX~Dissertation Template of \\Harbin Engineering University~(Version \version)}}



% 封面底部的中英文日期

\newcommand{\chinesethesistime}{2026~年~6~月}
\newcommand{\englishthesistime}{June, 2026}



%% 中文封面论文标题, 学科学位, 学科, 院系, 姓名, 导师姓名, 日期

\ctitle{水下绳驱机械臂结构设计及研制}  % 论文标题, 可手动断行
\cdegree{\cxueke\cxuewei}  % \cxueke 在 setup/Definition.tex 中定义, \cxuewei 在 setup/type.tex 中定义
\csubject{机械设计制造及其自动化}  % 学科
\caffil{机电工程学院}  % 院系
\cauthor{倪洪运}  % 或{某~~~~某}或超过三个字亦可
\csupervisor{刘星~~副教授}  % 导师名字
%\cassosupervisor{某~~~~某~~~~教授}  % 副导师, 如无可以不列此项
%\ccosupervisor{某某某~~~~教授~}  % 联合培养导师, 如无可以不列此项
\cdate{\chinesethesistime}  % 中文日期



%% 英文封面论文标题, 学科学位, 学科, 院系, 姓名, 导师姓名, 日期

\etitle{\englishthesistitle}  % 英文标题
\edegree{\exuewei \ of \exueke}  % \exueke 在 setup/Definition.tex 中定义, \exuewei 在 setup/type.tex 中定义
\esubject{Computer\hfill Science\hfill and\hfill Technology}  % 英文二级学科名
\eaffil{College\hfill of\hfill Computer\hfill Science\newline and Technology}  % 英文院系名, 如需换行请用 \newline, 不要用\\
\eauthor{XXX}  % 英文作者姓名
\esupervisor{Professor XXX}  % 英文导师姓名
%\ecosupervisor{Professor YYY}
%\eassosupervisor{Professor ZZZ}
\edate{\englishthesistime}  % 英文日期



%% 图书分类号和密级

\natclassifiedindex{TP309}  % 中图分类号
\internatclassifiedindex{681.324}  % 国际图书分类号
\statesecrets{公开} % 或{秘密}



%% 摘要内容

%\iffalse
%\BiAppendixChapter{摘~~~~要}{}  % 使用 winedt 编辑时, 为了在文档结构图(toc)中显示摘要, 需增加此句
%\fi



% 中文摘要
\cabstract{
  随着海洋探测技术的不断发展，微小型水下运载器对轻量化、紧凑化及低扰动作业机械臂的需求日益迫切。传统刚性水下机械臂存在自身质量大、运动惯性显著以及对运载器姿态扰动强等问题。为此，本文设计并研制了一款具有高负载自重比的三自由度水下绳驱机械臂，并对其结构、运动学、视觉定位及控制系统进行了系统深入的研究。

  首先，提出一种将驱动源集中后置于基座独立密封舱的绳驱总体方案，
  极大降低前端运动部件的质量与水下流体扰动。
  针对机械臂在运载器底部的完全折叠收纳需求，
  创新设计关联双铰链纯滚动肘部机构；
  通过滑轮组增力机构与单轴差动缠绕的导向走线方案，
  在不增加电机功率的前提下显著提升机械臂的承载能力。
  同时，基于极限受力工况，对核心轴类与连接件进行理论静力学计算与 ANSYS 有限元仿真，
  验证了本体结构的高可靠性。
  
  其次，开展绳驱机械臂的运动学建模与轨迹规划研究。
  基于改进 D-H 法建立了机械臂正运动学模型，
  采用牛顿-拉夫逊数值迭代法结合驱动-关节空间映射关系，
  成功解算了复杂的逆运动学。在关节空间与笛卡尔空间内，
  分别设计了平滑的梯形/多项式速度规划与直线/圆弧插补算法，
  通过 MATLAB 仿真验证所提算法能够有效抑制加速度突变，
  减小运动对水体的冲击。
  
  再次，构建基于单目视觉的目标定位与自主抓取策略。
  建立水下相机的针孔成像与畸变校正模型，
  基于 ArUco 特征标记与“眼在手外”构型，通过 
  Tsai-Lenz 法完成了高精度的手眼标定，
  得到从二维图像到机械臂物理基坐标系的空间位姿的映射。
  
  最后，完成试验样机的系统集成与性能测试。
  搭建以“树莓派 + STM32”异构主从架构的硬件平台，开发了基于 ROS 2 框架的分布式软件控制系统。
  样机实测最大下探深度达 730mm，
  具备 1.0kg 的额定有效负载能力。基础运动与视觉闭环抓取实验表明，
  该机械臂具备优异的运动响应性能与定位精度，验证整套系统的物理可行性与理论模型的正确性。
}

% 关键词
\ckeywords{水下机械臂；绳索驱动；结构设计；运动学建模；视觉抓取}



% 英文摘要
\eabstract{
  With the continuous advancement of ocean exploration technology, micro and small underwater vehicles increasingly demand lightweight, compact, and low-disturbance manipulators. Traditional rigid underwater manipulators suffer from large self-weight, significant motion inertia, and strong disturbances to the vehicle attitude. To address these problems, this thesis designs and develops a three-degree-of-freedom underwater cable-driven manipulator with a high load-to-weight ratio, and conducts systematic and in-depth research on its structure, kinematics, visual localization, and control system.

  First, a cable-driven overall scheme is proposed, in which all driving sources are centrally placed in an independent sealed chamber at the base, greatly reducing the mass of the front moving parts and underwater fluid disturbance. To meet the requirement of fully folding and stowing the manipulator beneath the vehicle, an innovative double-hinge pure rolling elbow mechanism is designed. By adopting a pulley-block force-amplifying mechanism and a single-axis differential winding routing scheme, the manipulator's payload capacity is significantly improved without increasing motor power. Meanwhile, based on extreme loading conditions, theoretical static calculations and ANSYS finite element simulations are carried out for key shafts and connectors, verifying the high reliability of the mechanical structure.
  
  Second, kinematic modeling and trajectory planning of the cable-driven manipulator are investigated. The forward kinematics model is established using the modified D-H method, and the complex inverse kinematics is successfully solved by combining the Newton–Raphson numerical iteration method with the mapping relationship between the driving space and joint space. In both joint space and Cartesian space, smooth trapezoidal/polynomial velocity profiles and linear/circular interpolation algorithms are designed. MATLAB simulations verify that the proposed algorithms can effectively suppress acceleration jumps and reduce hydrodynamic impact during motion.
  
  Third, a target localization and autonomous grasping strategy based on monocular vision is constructed. The pinhole imaging model and distortion correction model of the underwater camera are established. Using ArUco fiducial markers and the "eye-to-hand" configuration, high-precision hand–eye calibration is accomplished via the Tsai–Lenz method, yielding the spatial pose mapping from a 2D image to the manipulator's physical base coordinate frame.
  
  Finally, system integration and performance testing of the prototype are completed. A hardware platform based on a heterogeneous "Raspberry Pi + STM32" master–slave architecture is built, and a distributed software control system is developed under the ROS 2 framework. The prototype achieves a maximum downward reach of 730 mm and a rated payload capacity of 1.0 kg. Basic motion and visual closed-loop grasping experiments demonstrate that the manipulator possesses excellent motion response performance and positioning accuracy, validating the physical feasibility of the entire system and the correctness of the theoretical models.
  }

% 关键词
\ekeywords{Underwater Manipulator; Cable-Driven; Structural Design; Kinematics Modeling; Visual Grasping}



% 主要符号表

\NotationList{
  \begin{tabular}{ll}
    A & a matrix\\
    B & 登高\\
  \end{tabular}
}



% 生成封面

\makecover
\clearpage
